% ==============================================================================
% Capítulo 3: Trabalhos relacionados
% ==============================================================================
\chapter{Trabalhos relacionados}
\label{cap:relacionados}

\begin{instrucao}
\orientacao{O que vai em Trabalhos Relacionados}
Este capítulo mostra o estado da arte, ou seja, o que já foi feito sobre o
problema e onde o seu trabalho se encaixa. Ele responde a duas perguntas do
avaliador: ``o autor conhece a literatura do tema?'' e ``o que este trabalho
tem de diferente?''.

Evite a estrutura ``um parágrafo por artigo'', que vira uma lista de resumos.
Agrupe os trabalhos por abordagem (métodos exatos, heurísticas construtivas,
metaheurísticas, híbridos), compare-os entre si e aponte limitações. Termine
com um quadro comparativo e um parágrafo que posicione o seu trabalho em
relação aos demais.

Para encontrar trabalhos, use bases como Google Acadêmico, Scopus, IEEE
Xplore, ACM Digital Library e o Portal de Periódicos da Capes (acessível pela
rede da Ufes). Registre a busca (termos, bases e período) se ela fizer parte do
método.
\end{instrucao}

Os métodos propostos para o \acs{mkp} podem ser divididos em exatos,
heurísticos e híbridos \cite{freville2004multidimensional}. Este capítulo
discute os trabalhos mais relevantes de cada grupo para a abordagem proposta.

% ------------------------------------------------------------------------------
\section{Métodos exatos}
\label{sec:relacionados:exatos}

Os métodos exatos para o \acs{mkp} se baseiam, em sua maioria, em
\emph{branch-and-bound} com limitantes obtidos por relaxação linear,
lagrangiana ou substituta \cite{freville2004multidimensional}.
\textcite{puchinger2010multidimensional} mostraram que, em boa parte das
instâncias, os itens da solução ótima coincidem com os de maior eficiência na
relaxação linear, exceto por um pequeno conjunto de itens intermediários,
chamado de núcleo. Resolver o problema restrito ao núcleo reduz o tempo de
solução, mas instâncias com muitas restrições continuam fora do alcance dos
métodos exatos em tempo limitado.

% ------------------------------------------------------------------------------
\section{Métodos heurísticos e híbridos}
\label{sec:relacionados:heuristicos}

Entre as metaheurísticas, o \ac{ag} de \textcite{chu1998genetic} tornou-se uma
referência para o problema. O método usa um operador de reparo que remove e
insere itens segundo uma razão de utilidade derivada da relaxação do problema,
o que garante que toda solução gerada seja viável. Os autores também
publicaram as 270 instâncias usadas nos experimentos, que se tornaram o
conjunto padrão de comparação na área.

Abordagens híbridas combinam programação matemática e busca heurística.
\textcite{vasquez2001hybrid} usam a relaxação linear para fixar o número de
itens da solução e aplicam busca tabu nas regiões indicadas por ela, obtendo
alguns dos melhores resultados conhecidos para as maiores instâncias da
OR-Library. Em contrapartida, o método exige resolver vários programas lineares
e tem diversos parâmetros, o que dificulta sua reprodução. A combinação de
núcleo com algoritmo memético e resolução exata proposta por
\textcite{puchinger2010multidimensional} segue a mesma linha.

% ------------------------------------------------------------------------------
\section{Síntese e posicionamento do trabalho}
\label{sec:relacionados:sintese}

O \cref{qua:relacionados} resume as abordagens discutidas e as compara com a
proposta deste trabalho. Observa-se que os métodos de melhor desempenho
dependem da resolução repetida de relaxações lineares ou de muitos parâmetros.
A heurística proposta procura um meio-termo, com poucos parâmetros, sem
resolver programas lineares durante a busca e com desempenho competitivo em
tempo limitado.

\begin{quadro}[htbp]
  \centering
  \caption{Comparação entre os trabalhos relacionados e a proposta}
  \label{qua:relacionados}
  \small
  \begin{tabularx}{\textwidth}{|l|X|c|}
    \hline
    \textbf{Trabalho} & \textbf{Abordagem} & \textbf{Tipo} \\
    \hline
    \textcite{chu1998genetic} & \Acs{ag} com operador de reparo baseado em
      razão de utilidade & Heurístico \\
    \hline
    \textcite{vasquez2001hybrid} & Relaxação linear para guiar uma busca tabu
      & Híbrido \\
    \hline
    \textcite{puchinger2010multidimensional} & Redução ao núcleo do problema
      com algoritmo memético e resolução exata & Híbrido \\
    \hline
    Este trabalho & \Acs{ils} com perturbação adaptativa e solução inicial gulosa
      & Heurístico \\
    \hline
  \end{tabularx}
  \fonteautor
\end{quadro}

\begin{instrucao}
\orientacao{Quadro ou tabela?}
A NBR 14724 distingue os dois. O \textbf{quadro} organiza informações
textuais em linhas e colunas e tem bordas fechadas, com linhas verticais. A
\textbf{tabela} apresenta dados, em geral numéricos, e segue as normas do
IBGE, com bordas laterais abertas e sem linhas verticais (veja a
\cref{tab:resultados} no \cref{cap:experimentos}). Quadros e tabelas têm listas
e numeração próprias. A identificação fica acima e a fonte abaixo, como nas
figuras.
\begin{exemplocodigo}
\begin{quadro}[htbp]
  \centering
  \caption{Comparação entre os trabalhos}
  \label{qua:relacionados}
  \small
  \begin{tabularx}{\textwidth}{|l|X|c|}
    \hline
    \textbf{Trabalho} & \textbf{Abordagem} & \textbf{Tipo} \\
    \hline
    \textcite{chu1998genetic} & Algoritmo genético ... & Heurístico \\
    \hline
  \end{tabularx}
  \fonteautor
\end{quadro}
\end{exemplocodigo}
A coluna \verb|X| do pacote \verb|tabularx| ocupa a largura que sobra e quebra
o texto em várias linhas. Para montar tabelas e quadros mais elaborados, o site
\href{https://www.tablesgenerator.com/}{Tables Generator}
(\url{https://www.tablesgenerator.com/}) gera o código LaTeX a partir de uma
planilha. Depois de colar o código, ajuste-o ao padrão do modelo (ambiente
\verb|table| ou \verb|quadro|, \verb|\caption| acima e \verb|\fonte| abaixo).
\end{instrucao}
